\paragraph{Cited theorem.}
Routley and Meyer \cite[pp.~53--73]{RoutleyMeyer1972} pair the entailment
systems specified in that source with ternary relational semantics and prove
the corresponding source-scoped semantic results.  The controlling hypotheses
are the selected source calculus and its matching postulates on the ternary
relation, normal points, and star operation; changing either side changes the
result.  This citation supports the displayed semantic architecture and the
need for a matched pair.  It is not a completeness theorem for the deliberately
broad base frames defined in this chapter.

\paragraph{Internal proof boundary.}
A local completeness theorem would require selecting one existing system and
reproducing its exact formula language; every axiom and inference rule; the
context, label, and structural policies; all side conditions; finite derivation
objects; and the precise corresponding frame postulates.  A complete proof
would then need rule-by-rule soundness and a canonical-model or
saturated-branch construction with an existence lemma and truth lemma for
every constructor.  Those source-dependent data are not supplied.  The chapter
therefore uses only the explicitly displayed satisfaction clauses and
normal-world validity, and introduces no calculus name or
derivability--validity biconditional.

